\documentclass[12pt,a4paper]{article}
\usepackage[utf8]{inputenc}
\usepackage[T5]{fontenc}
\usepackage{amsmath, amssymb}
\usepackage{graphicx}
\usepackage{ifthen}

\newboolean{showsolutions}
\setboolean{showsolutions}{true}

% --- ĐỊNH NGHĨA CÁC LỆNH ẢO CHO CÔNG CỤ LATEX2JSON CỦA WEBSITE ---
\newcommand{\doctitle}[1]{\begin{center}\Large\textbf{#1}\end{center}}
\newcommand{\docdesc}[1]{\textbf{Mô tả:} #1}
\newcommand{\docgrade}[1]{\textbf{Lớp:} #1}
\newcommand{\docstatus}[1]{\textbf{Trạng thái:} #1}
\newcommand{\doctype}[1]{\textbf{Loại:} #1}
\newcommand{\doctopics}[1]{\textbf{Chủ đề:} #1}

\newenvironment{textblock}{\vspace{1em}}{}

\newenvironment{lesson}[2]{
	\vspace{1em}\noindent\textbf{\Large #1}\\
	\textit{#2}\vspace{0.5em}
}{}

\newcommand{\image}[3]{
	\begin{center}
		\includegraphics[width=#1\textwidth]{#3} \\
		\textit{#2}
	\end{center}
}

\newenvironment{quiz}[1]{
	\vspace{1em}\noindent\textbf{\Large Kiểm tra: #1}
}{}

\newenvironment{mcq}[4]{
	\vspace{1em}\noindent\textbf{Câu hỏi:} #1 \\
	\ifthenelse{\boolean{showsolutions}}{
		\textit{(Đáp án đúng: #2, Điểm: #3) - Giải thích: #4}
	}{}
	\begin{itemize}
	}{
	\end{itemize}
}
\newcommand{\option}[1]{\item #1}

\newenvironment{truefalse}[3]{
	\vspace{1em}\noindent\textbf{Đúng/Sai:} #1 \\
	\ifthenelse{\boolean{showsolutions}}{
		\textit{(Điểm: #2) - Giải thích: #3}
	}{}
	\begin{itemize}
	}{
	\end{itemize}
}
\newcommand{\statement}[2]{\item [\textbf{#1}] #2}

\newcommand{\shortanswer}[4]{
    \vspace{1em}\noindent\textbf{Trả lời ngắn (Điểm: #3):}

    \noindent #1

	\ifthenelse{\boolean{showsolutions}}{
		\vspace{0.5em}
		\noindent\textit{Đáp án: #2 - Giải thích:}
		\noindent #4
	}{}
}

\newcommand{\essay}[3]{
    \vspace{1em}\noindent\textbf{Tự luận (Điểm: #2):}
    
    \noindent #1
    
	\ifthenelse{\boolean{showsolutions}}{
		\vspace{0.5em}
		\noindent\textbf{Gợi ý / Đáp án:}
		\noindent #3
	}{}
}

\begin{document}

\doctitle{ĐỀ THI RÈN LUYỆN TÍCH PHÂN - ĐỀ 01}
\docdesc{Đề kiểm tra rèn luyện chuyên đề Tích phân - Đề số 1}
\docgrade{12}
\docstatus{draft}
\doctype{test}
\doctopics{Giải tích, Tích phân}

% =========================================================================
% PHẦN I: CÂU TRẮC NGHIỆM NHIỀU PHƯƠNG ÁN LỰA CHỌN
% =========================================================================

\begin{quiz}{Phần 1. Câu trắc nghiệm nhiều phương án lựa chọn}

\begin{mcq}{Gọi $S$ là diện tích hình phẳng giới hạn bởi đồ thị hàm số $y = f(x)$, trục hoành và hai đường thẳng $x = a, x = b$ (phần gạch chéo ở hình vẽ bên dưới). Khẳng định nào dưới đây là sai? \image{0.45}{}{images_ChuDe02_TichPhan/d1_c1.png}}{1}{1}{Chọn B.}
	\option{$S = \int_a^b |f(x)|\,\mathrm{d}x$.}
	\option{$S = -\int_a^b f(x)\,\mathrm{d}x$.}
	\option{$S = \left|\int_a^b f(x)\,\mathrm{d}x\right|$.}
	\option{$S = \int_a^b f(x)\,\mathrm{d}x$.}
\end{mcq}

\begin{mcq}{Diện tích hình thang cong giới hạn bởi đồ thị hàm số $y = \sin x$, trục hoành và hai đường thẳng $x = \frac{\pi}{6}, x = \frac{5\pi}{6}$ là \image{0.45}{}{images_ChuDe02_TichPhan/d1_c2.png}}{3}{1}{Chọn D.}
	\option{$\frac{1}{2} - \frac{\sqrt{3}}{2}$.}
	\option{$\frac{\sqrt{3}}{2} - \frac{1}{2}$.}
	\option{$0$.}
	\option{$\sqrt{3}$.}
\end{mcq}

\begin{mcq}{Biết $\int_1^3 f(x)\,\mathrm{d}x = 3$. Giá trị của $\int_1^3 2f(x)\,\mathrm{d}x$ bằng}{2}{1}{Chọn C.}
	\option{$5$.}
	\option{$9$.}
	\option{$6$.}
	\option{$\frac{3}{2}$.}
\end{mcq}

\begin{mcq}{Tính $I = \int_0^1 \left(\frac{1}{2x+1} + 3\sqrt{x}\right)\mathrm{d}x$.}{2}{1}{Chọn C.}
	\option{$1 + \ln\sqrt{3}$.}
	\option{$2 + \ln 3$.}
	\option{$2 + \ln\sqrt{3}$.}
	\option{$4 + \ln 3$.}
\end{mcq}

\begin{mcq}{Giá trị của $\int_0^1 (2025x^{2024} - 1)\,\mathrm{d}x$ là}{0}{1}{Chọn A.}
	\option{$0$.}
	\option{$2^{2023} + 1$.}
	\option{$2^{2023} - 1$.}
	\option{$1$.}
\end{mcq}

\begin{mcq}{Cho hàm số $f(x)$ liên tục trên $\mathbb{R}$ và $\int_0^2 [f(x)+3x^2]\,\mathrm{d}x = 10$. Tính $\int_0^2 f(x)\,\mathrm{d}x$.}{3}{1}{Chọn D.}
	\option{$-18$.}
	\option{$-2$.}
	\option{$18$.}
	\option{$2$.}
\end{mcq}

\begin{mcq}{Nếu $\int_1^2 f(x)\,\mathrm{d}x = 3, \int_1^2 [3f(x)-g(x)]\,\mathrm{d}x = 2$ thì $\int_1^2 g(x)\,\mathrm{d}x$ bằng}{3}{1}{Chọn D.}
	\option{$11$.}
	\option{$5$.}
	\option{$1$.}
	\option{$7$.}
\end{mcq}

\begin{mcq}{Nếu $\int_0^2 f(x)\,\mathrm{d}x = 3, \int_0^2 g(x)\,\mathrm{d}x = -5$ thì $\int_0^2 [3f(x)+4g(x)]\,\mathrm{d}x$ bằng}{2}{1}{Chọn C.}
	\option{$29$.}
	\option{$-3$.}
	\option{$-11$.}
	\option{$4$.}
\end{mcq}

\begin{mcq}{Cho hai tích phân $\int_{-2}^5 f(x)\,\mathrm{d}x = 8, \int_{-2}^5 g(x)\,\mathrm{d}x = -3$. Tính $\int_{-2}^5 [f(x)-4g(x)-1]\,\mathrm{d}x$.}{1}{1}{Chọn B.}
	\option{$-11$.}
	\option{$13$.}
	\option{$27$.}
	\option{$3$.}
\end{mcq}

\begin{mcq}{Biết rằng, có hai giá trị của số thực $a$ là $a_1$ và $a_2$ ($0 < a_1 < a_2$) thỏa mãn $\int_1^a (2x-3)\,\mathrm{d}x = 0$. Hãy tính giá trị của biểu thức $T = 3^{a_1} + 3^{a_2} + \log_2\left(\frac{a_2}{a_1}\right)$.}{2}{1}{Chọn C.}
	\option{$T = 26$.}
	\option{$T = 12$.}
	\option{$T = 13$.}
	\option{$T = 28$.}
\end{mcq}

\begin{mcq}{Cho hàm số $y = f(x)$ đạt cực trị tại $x = -1$ và có đạo hàm cấp hai liên tục trên $\mathbb{R}$. Đồ thị hàm số $f(x)$ như hình vẽ bên dưới. Biết rằng đường thẳng $d$ là một tiếp tuyến của $(C)$ tại điểm có hoành độ $x = 3$. Tính $\int_{-1}^3 f''(x)\,\mathrm{d}x$. \image{0.45}{}{images_ChuDe02_TichPhan/d1_c11.png}}{3}{1}{Chọn D.}
	\option{$0$.}
	\option{$3$.}
	\option{$-2$.}
	\option{$1$.}
\end{mcq}

\begin{mcq}{Cho hàm số $y = f(x)$ liên tục, luôn dương trên $[0; 3]$ và thỏa mãn $\int_0^3 f(x)\,\mathrm{d}x = 4$. Khi đó giá trị của tích phân $K = \int_0^3 \left(e^{1+\ln f(x)} + 4\right)\mathrm{d}x$ là}{2}{1}{Chọn C.}
	\option{$14 + 3e$.}
	\option{$14 + 4e$.}
	\option{$12 + 4e$.}
	\option{$12 + 3e$.}
\end{mcq}

\begin{mcq}{Biết $I = \int_1^2 \frac{x^2+2x}{x+1}\,\mathrm{d}x = \frac{5}{a} + \ln b - \ln c$. Tính giá trị biểu thức $S = a - b + c$.}{1}{1}{Chọn B.}
	\option{$S = 7$.}
	\option{$S = 3$.}
	\option{$S = -3$.}
	\option{$S = 1$.}
\end{mcq}

\begin{mcq}{Biết $\int_1^2 \frac{\mathrm{d}x}{x\sqrt{x+2} + (x+2)\sqrt{x}} = \sqrt{a} + \sqrt{b} - c (a, b, c \in \mathbb{Z}^+)$. Tính $P = a + b + c$.}{1}{1}{Chọn B.}
	\option{$P = 2$.}
	\option{$P = 8$.}
	\option{$P = 46$.}
	\option{$P = 22$.}
\end{mcq}

\end{quiz}

% =========================================================================
% PHẦN II: CÂU TRẮC NGHIỆM ĐÚNG SAI
% =========================================================================

\begin{quiz}{Phần 2. Câu trắc nghiệm đúng sai}

\begin{truefalse}{Cho các hàm số $f(x)$ và $g(x)$ là hai hàm số liên tục trên $\mathbb{R}$. Xét tính đúng sai của các khẳng định sau:}{1}{Đáp án: a) Đúng, b) Sai, c) Sai, d) Đúng.}
	\statement{true}{Nếu $F(x)$ là một nguyên hàm của hàm số $f(x)$ thỏa mãn $\int_0^9 f(x)\,\mathrm{d}x = 9$ và $F(0) = 3$ thì $F(9) = 12$.}
	\statement{false}{Nếu $\int_1^3 [f(x)+3g(x)]\,\mathrm{d}x = 10$ và $\int_1^3 [2f(x)-g(x)]\,\mathrm{d}x = 6$ thì $\int_1^3 [f(x)+g(x)]\,\mathrm{d}x = 7$.}
	\statement{false}{$\int_a^b f(x)\,\mathrm{d}x - \int_a^c f(x)\,\mathrm{d}x = \int_b^c f(x)\,\mathrm{d}x$.}
	\statement{true}{Nếu $\int_1^2 f(x)\,\mathrm{d}x = 3$ và $\int_2^5 f(x)\,\mathrm{d}x = -1$ thì $\int_1^5 f(x)\,\mathrm{d}x = 2$.}
\end{truefalse}

\begin{truefalse}{Cho hàm số $f(x)$ liên tục trên $\mathbb{R}$. Xét tính đúng sai của các khẳng định sau:}{1}{Đáp án: a) Đúng, b) Sai, c) Đúng, d) Đúng.}
	\statement{true}{Với $a, b$ là các tham số thực. Giá trị của tích phân $\int_0^b (3x^2 - 2ax - 1)\,\mathrm{d}x$ bằng $b^3 - ab^2 - b$.}
	\statement{false}{Nếu $\int_0^2 [f(x)+3x^2]\,\mathrm{d}x = 10$ thì $\int_0^2 f(x)\,\mathrm{d}x = -2$.}
	\statement{true}{Nếu $\int_{-2}^1 [2f(x)-1]\,\mathrm{d}x = 3$ thì $\int_{-2}^1 f(x)\,\mathrm{d}x = 3$.}
	\statement{true}{Cho $I = \int_0^1 \frac{\mathrm{d}x}{\sqrt{2x+m}}, m$ là số thực dương. Nếu $I \ge 1$ thì $0 < m \le \frac{1}{4}$.}
\end{truefalse}

\begin{truefalse}{Một người đi xe máy đang di chuyển với tốc độ $54\text{ km/h}$, thì bất ngờ thấy chướng ngại vật cách người đó $50\text{ m}$. Sau đó một giây thì người này bắt đầu phanh gấp để xe chuyển động chậm dần đều với vận tốc $v(t) = -10t + 15$, trong đó $t$ là thời gian được tính bằng giây kể từ thời điểm bắt đầu phanh. Gọi $s(t)$ là quãng đường mà xe máy đi được trong $t$ giây. Xét tính đúng sai của các khẳng định sau:}{1}{Đáp án: a) Đúng, b) Đúng, c) Sai, d) Đúng.}
	\statement{true}{Quãng đường $s(t)$ xe máy đi được là một nguyên hàm của hàm số $v(t)$.}
	\statement{true}{$s(t) = -5t^2 + 15t$.}
	\statement{false}{Thời gian để ô tô dừng hẳn lúc $1{,}5\text{s}$ khi ô tô đi được $19\text{m}$ kể từ lúc bắt đầu phanh xe.}
	\statement{true}{Xe máy đã dừng trước và không va chạm vào chướng ngại vật.}
\end{truefalse}

\begin{truefalse}{Một vật chuyển động với vận tốc $10\text{ m/s}$ thì tăng tốc với gia tốc được tính theo thời gian là $a(t) = t^2 + 3t$. Xét tính đúng sai của các khẳng định sau:}{1}{Đáp án: a) Đúng, b) Sai, c) Đúng, d) Đúng.}
	\statement{true}{Vận tốc được biểu diễn theo thời gian là $v(t) = \frac{1}{3}t^3 + \frac{3}{2}t^2 + 10$.}
	\statement{false}{$v(1) = 12\text{ m/s}$.}
	\statement{true}{Quãng đường vật đi được biểu diễn theo thời gian là $s(t) = \frac{t^4}{12} + \frac{t^3}{2} + 10t$.}
	\statement{true}{Quãng đường vật đi được trong khoảng thời gian $3\text{s}$ kể từ khi vật bắt đầu tăng tốc là $50{,}25\text{ (m)}$.}
\end{truefalse}

\end{quiz}

% =========================================================================
% PHẦN III: CÂU TRẮC NGHIỆM TRẢ LỜI NGẮN
% =========================================================================

\begin{quiz}{Phần 3. Câu trắc nghiệm trả lời ngắn}

\shortanswer{Một vật chuyển động chậm dần với vận tốc $v(t) = 160 - 10t\text{ (m/s)}$. Quãng đường $S$ mà vật di chuyển trong khoảng thời gian từ thời điểm $t = 0\text{ (s)}$ đến thời điểm vật dừng lại là bao nhiêu?}{$1280$}{1}{Đáp án: 1280.}

\shortanswer{Cho hàm số $y = f(x)$ xác định, liên tục và có nguyên hàm trên $[-2; 4]$ đồng thời có đồ thị như hình vẽ bên. Tích phân $I = \int_{-2}^4 f(x)\,\mathrm{d}x$ bằng bao nhiêu? \image{0.45}{}{images_ChuDe02_TichPhan/sa_c2.png}}{$4$}{1}{Đáp án: 4.}

\shortanswer{Cho hàm số $y = f(x)$ là hàm đa thức bậc ba có đồ thị như hình vẽ. Tích phân $\int_0^3 f(x)\,\mathrm{d}x$ bằng bao nhiêu? (làm tròn đến 2 chữ số sau dấu thập phân). \image{0.45}{}{images_ChuDe02_TichPhan/sa_c3.png}}{$2{,}25$}{1}{Đáp án: 2,25.}

\shortanswer{Một mol khí $\text{N}_2$ thực hiện quá trình giãn nở đẳng áp như hình bên từ trạng thái (1) sang trạng thái (2). Biết công thực hiện được tính bằng công thức $A = \int_1^2 p\,\mathrm{d}v$ và $1\text{Pa.m}^3 = 1\text{J}$. Độ lớn công thực hiện ở quá trình giãn nở đẳng áp từ trạng thái (1) sang trạng thái (2) là bao nhiêu? \image{0.45}{}{images_ChuDe02_TichPhan/sa_c4.png}}{$7500$}{1}{Đáp án: 7500.}

\shortanswer{Cho tích phân $I = \int_a^b |x+1|\,\mathrm{d}x = ma^2 + nb^2 + pa + qb$ với $a, b, m, n, p, q \in \mathbb{R}$ và $a < -1 < b$. Giá trị của biểu thức $T = e^{mn} \cdot \ln(pq)$ bằng bao nhiêu?}{$0$}{1}{Đáp án: 0.}

\shortanswer{Cho hàm số $f(x) = 6x^2 + \int_0^2 f(x)\,\mathrm{d}x$. Giá trị của $f(3)$ bằng bao nhiêu?}{$38$}{1}{Đáp án: 38.}

\shortanswer{Cho hàm số $y = f(x)$ liên tục trên $\mathbb{R}$ thỏa mãn $\int_0^3 [x^2 f'(x) + 2x f(x)]\,\mathrm{d}x = 18$. Giá trị của $f(3)$ bằng bao nhiêu?}{$2$}{1}{Đáp án: 2.}

\shortanswer{Cho hàm số $y = f(x)$ liên tục trên $\mathbb{R}$ và có đạo hàm cấp hai liên tục trên $\mathbb{R}$. Biết hàm số $y = f(x)$ đạt cực trị tại $x = -1$, có đồ thị như hình vẽ và đường thẳng $\Delta$ là tiếp tuyến của đồ thị hàm số tại điểm có hoành độ bằng 2. Tích phân $\int_1^4 f''(x-2)\,\mathrm{d}x$ bằng bao nhiêu? \image{0.45}{}{images_ChuDe02_TichPhan/sa_c8.png}}{$2$}{1}{Đáp án: 2.}

\shortanswer{Cho hàm số $f(x)$ liên tục, có đạo hàm liên tục trên $\mathbb{R}$ và có đồ thị như hình vẽ. Tích phân $I = \int_0^1 f'(5x-3)\,\mathrm{d}x$ bằng bao nhiêu? \image{0.45}{}{images_ChuDe02_TichPhan/sa_c9.png}}{$-0{,}6$}{1}{Đáp án: -0,6.}

\end{quiz}

\end{document}
